\section{Considerações finais}
\label{sec:consideracoes}

Esta pesquisa busca contribuir para a compreensão das relações entre exposição ocupacional à inteligência artificial e transições no mercado de trabalho brasileiro. Sua contribuição pretendida está na organização de uma análise baseada na PNAD Contínua, com documentação das escolhas metodológicas e distinção entre exposição potencial, adoção efetiva e resultados observados. A avaliação dessa contribuição dependerá da consistência dos procedimentos e das evidências apresentadas.

As conclusões da dissertação deverão retomar a pergunta de pesquisa e discutir o que os resultados permitem afirmar, preservando os limites do desenho adotado. A ausência de uma medida direta de adoção da inteligência artificial e as possíveis mudanças na coleta e na composição da amostra deverão ser consideradas. Pesquisas futuras poderão explorar outras medidas de exposição, diferentes recortes e estratégias adicionais de validação.

Este capítulo é provisório e integra o exercício ia-6.1. Não estabelece conclusões empíricas novas nem antecipa resultados que ainda não tenham sido verificados e discutidos pelo autor.